\section*{Zusammenfassung}
\addcontentsline{toc}{section}{Zusammenfassung}

Beim Jassen werden die Karten am Ende einer Runde von Hand gezählt. Das
kostet Zeit und ist besonders für ungeübte Spielende mühsam. Ziel dieser
Arbeit ist ein Modell, das in einem Videostream die oberste Karte eines
Stapels in Echtzeit erkennt und so das Zählen mit dem Smartphone ermöglicht.

Das Training erfolgt ausschliesslich mit synthetisch erzeugten Bildern. Ein
Generator setzt aus 72 fotografierten Jasskarten Stapel mit unterschiedlichen
Kameraperspektiven, Beleuchtungen und Tischoberflächen zusammen und erzeugt
die Labels automatisch. Drei Ansätze auf Basis von YOLO11 werden verglichen:

\begin{itemize}
  \item Variante~A klassifiziert das gesamte Bild.
  \item Variante~B lokalisiert zuerst die oberste Karte und klassifiziert
    anschliessend ihren entzerrten Ausschnitt.
  \item Variante~C erkennt Klasse und Position der obersten Karte mit einer
    achsenparallelen Bounding Box in einem Schritt.
\end{itemize}

Geprüft werden die Modelle auf eigenen, von Hand gelabelten realen Aufnahmen.
Die letzte Datenmengenreihe vergleicht sechs Bildanzahlen auf demselben
Validierungsset mit 813 Frames. Variante~C verbessert sich mit mehr
Trainingsbildern und erreicht bei 200'000 Bildern eine \map{50} von 0.908.
Ein zusätzlicher Lauf mit 500'000 Bildern bringt keinen weiteren Gewinn.
Sie liefert Klasse und Position mit einem Modell und läuft in den Apps
in Echtzeit. Variante~A liefert keine Position. Bei Variante~B ist die
gesamte zweistufige Erkennung auf diesem Stand nicht verlässlich gemessen.

Dies aufgrund der bereits in der ersten Stufe von Variante~B schlechteren Ergebnisse gegenüber C, daraus würden auch für die Stufe B2 keine Verbsserungen erfolgen.

Die Ergebnisse sprechen deshalb am stärksten für Variante~C.

Das Modell wurde in eine iPhone- und eine Android-App zur Kartenzählung
integriert. Fehlvorhersagen, insbesondere bei Karten im Wurf, bleiben eine
Herausforderung. Eine Stabilitätsregel kann kurze Fehlerfolgen abfangen,
behebt aber nicht alle Fehlalarme. Da viele Validierungsbilder aus denselben
Videos stammen und ein unabhängiges Testset fehlt, sind die gemessenen
Trefferquoten vorsichtig einzuordnen.


Das Training läuft automatisiert auf
zeitweise gemieteten Cloud-GPUs und ist über Commit, Seed und Konfiguration
nachvollziehbar.

\clearpage
